\documentclass[a4paper]{article}
\usepackage[a4paper,margin=0mm]{geometry}
\usepackage{fontspec}
\usepackage{unicode-math}
\usepackage{tikz}
\usepackage[table]{xcolor}
\usepackage[absolute]{textpos}
\usepackage{array}

\setmainfont{Libertinus Serif}
\setsansfont{Libertinus Sans}
\setmathfont{Libertinus Math}

\definecolor{accent}{HTML}{0F6A52}
\definecolor{badge}{HTML}{A9CFC2}
\definecolor{ink}{HTML}{2A2E30}

\setlength{\TPHorizModule}{1mm}
\setlength{\TPVertModule}{1mm}
\setlength{\parindent}{0pt}
\pagestyle{empty}


\newcommand{\pagefooter}[1]{%
\begin{textblock*}{194mm}(8mm,286.41mm)
  \begin{tikzpicture}\fill[accent] (0,0) rectangle (194mm,0.6mm);\end{tikzpicture}
\end{textblock*}
\begin{textblock*}{63.59mm}(8mm,288.27mm)
  {\sffamily\bfseries\color{accent}\fontsize{13}{16}\selectfont SAT 1200+}
\end{textblock*}
\begin{textblock*}{10mm}(188mm,288.27mm)
  \centering{\sffamily\color{accent}\fontsize{14}{16}\selectfont #1}
\end{textblock*}}
\newcommand{\frontheader}[1]{%
\begin{textblock*}{194mm}(16mm,16mm)
\begin{tikzpicture}\node[font=\sffamily\bfseries\fontsize{15}{16}\selectfont,text=ink,anchor=west] at (0,0) {#1};\end{tikzpicture}
\end{textblock*}
\begin{textblock*}{194mm}(16mm,22.5mm)
\begin{tikzpicture}\fill[accent] (0,0) rectangle (194mm,0.5mm);\end{tikzpicture}
\end{textblock*}}

\begin{document}
\frontheader{Mock Tests}
\pagefooter{@@PAGE@@}
\mbox{}
\begin{textblock*}{170mm}(20mm,40mm)
{\fontsize{11}{17}\selectfont\color{ink}
\textbf{10 ta to'liq mock test.} Har bir mock --- haqiqiy Digital SAT Math tuzilishi: ikki modul, har biri 22 savol, 35 daqiqa.

\vspace{4mm}
\textbf{1-modul} --- hamma uchun bir xil, oson, o'rta va qiyin savollar aralashmasi (taxminan 8 oson, 11 o'rta, 3 qiyin). Haqiqiy testda 1-moduldagi natijangiz 2-modulning qaysi versiyasini olishingizni belgilaydi.

\vspace{2mm}
\textbf{2-modul} --- bu kitobda doim \emph{qiyinroq yo'nalish} (taxminan 2 oson, 10 o'rta, 10 qiyin). Chunki yuqori ballga faqat shu yo'l olib boradi.

\vspace{4mm}
\textbf{Vaqt rejasi (har bir modulda):}

\vspace{1mm}
\par\vspace{2.2mm}\noindent\hangindent=7mm\hangafter=1 \hspace{1mm}{\color{accent}\textbullet}\hspace{3.5mm}1--15-savollar: taxminan 18 daqiqa --- tez va aniq, shubhali savolni belgilab keting.
\par\vspace{2.2mm}\noindent\hangindent=7mm\hangafter=1 \hspace{1mm}{\color{accent}\textbullet}\hspace{3.5mm}16--22-savollar: taxminan 17 daqiqa --- eng qiyin savollar oxirida joylashgan.
\par\vspace{2.2mm}\noindent\hangindent=7mm\hangafter=1 \hspace{1mm}{\color{accent}\textbullet}\hspace{3.5mm}So'nggi 2--3 daqiqa: belgilangan savollarga qaytish va barcha bo'sh javoblarni to'ldirish --- bo'sh javob nol ball, taxmin esa yo'q.\par

\vspace{4mm}
\textbf{Qanday ishlash kerak:} har bir modulni o'z vaqtida, kalkulyator bilan (Desmos) bajaring. Modullar orasida 5--10 daqiqa dam oling. Javoblarni \textbf{tekshirmang} --- ikkala modulni tugatgach, kitob oxiridagi \textbf{Answer sheets} bo'limida (Mock Test raqami bo'yicha) javoblarni solishtiring. Har bir xato savol uchun uning sohasi (ALG, ADV, PSDA, GT) va darajasini belgilab boring --- shu ro'yxat sizga qaysi darsga qaytish kerakligini aytadi.

\vspace{4mm}
\textit{Haftasiga 1--2 ta mock yetarli: hammasini bir haftada ishlash shoshqaloqlikka va noto'g'ri xulosalarga olib keladi.}
}
\end{textblock*}
\end{document}
